\documentclass[11pt,letterpaper]{article}
% ==============================================================================
% Shared Preamble: Principles of Economics Modular Curriculum
% Core Colors: SlateBlue (#1E293B) & Terracotta (#C2410C)
% ==============================================================================

\usepackage[utf8]{inputenc}
\usepackage[T1]{fontenc}
\usepackage{lmodern}
\usepackage[margin=1in]{geometry}
\usepackage{amsmath,amssymb,amsfonts,amsthm}
\usepackage{booktabs}
\usepackage{microtype}
\usepackage{enumitem}
\usepackage{fancyhdr}
\usepackage{titlesec}
\usepackage{xcolor}
\usepackage{tikz}
\usepackage{pgfplots}
\usepackage{mdframed}
\usepackage[hidelinks,colorlinks=true,linkcolor=SlateBlue,citecolor=SlateBlue,urlcolor=Terracotta]{hyperref}

\pgfplotsset{compat=1.18}
\usetikzlibrary{arrows.meta,positioning,calc,patterns,decorations.pathreplacing}

% --- Color Palette Definitions ---
\definecolor{SlateBlue}{HTML}{1E293B}      % Primary Accent / Structure
\definecolor{Terracotta}{HTML}{C2410C}     % Secondary Accent / Emphasis
\definecolor{ForestGreen}{HTML}{15803D}    % Equilibrium / Positive Shocks
\definecolor{LightSlate}{HTML}{F8FAFC}     % Box Background Fill
\definecolor{BorderGray}{HTML}{CBD5E1}     % Neutral Borders
\definecolor{LightTerracotta}{HTML}{FFF7ED}% Cautionary / Principle Fill

% --- Section Heading Styling ---
\titleformat{\section}
  {\normalfont\Large\bfseries\color{SlateBlue}}{\thesection}{1em}{}[\titlerule]
\titleformat{\subsection}
  {\normalfont\large\bfseries\color{SlateBlue}}{\thesubsection}{1em}{}
\titleformat{\subsubsection}
  {\normalfont\normalsize\bfseries\color{Terracotta}}{\thesubsubsection}{1em}{}

% --- Header and Footer Configuration ---
\pagestyle{fancy}
\fancyhf{}
\renewcommand{\headrulewidth}{0.4pt}
\renewcommand{\footrulewidth}{0.4pt}
\fancyhead[L]{\small\scshape\color{SlateBlue}\leftmark}
\fancyhead[R]{\small\color{SlateBlue}Principles of Economics}
\fancyfoot[C]{\thepage}

% --- Custom Environments (mdframed) ---

% 1. Core Principle Box
\newmdenv[
  linecolor=Terracotta,
  linewidth=1.5pt,
  backgroundcolor=LightTerracotta,
  innertopmargin=10pt,
  innerbottommargin=10pt,
  innerrightmargin=12pt,
  innerleftmargin=12pt,
  roundcorner=4pt,
  skipabove=12pt,
  skipbelow=12pt
]{coreprinciple}

% 2. Empirical Case Study Box
\newmdenv[
  linecolor=SlateBlue,
  linewidth=1.5pt,
  backgroundcolor=LightSlate,
  innertopmargin=10pt,
  innerbottommargin=10pt,
  innerrightmargin=12pt,
  innerleftmargin=12pt,
  roundcorner=4pt,
  skipabove=12pt,
  skipbelow=12pt
]{casestudy}

% 3. Key Takeaways Box
\newmdenv[
  linecolor=BorderGray,
  linewidth=1pt,
  backgroundcolor=white,
  innertopmargin=10pt,
  innerbottommargin=10pt,
  innerrightmargin=12pt,
  innerleftmargin=12pt,
  roundcorner=4pt,
  skipabove=12pt,
  skipbelow=12pt
]{takeaways}

% Environment Title Helpers
\newcommand{\boxheader}[2]{{\noindent\bfseries\color{#1}\Large #2}\vspace{0.4em}\par}
\usepackage{tikz}
\usepackage{pgfplots}
\pgfplotsset{compat=1.18}
\usetikzlibrary{arrows.meta,positioning,calc,shapes.geometric,patterns,decorations.pathreplacing,intersections}

\tikzset{
  econ axis/.style={->, >=Stealth, thick, color=SlateBlue!90!black},
  econ grid/.style={thin, dashed, color=gray!30},
  econ guide/.style={dashed, color=gray!60, thin},
  demand curve/.style={very thick, color=SlateBlue},
  supply curve/.style={very thick, color=Terracotta},
  mb curve/.style={very thick, color=SlateBlue},
  mc curve/.style={very thick, color=Terracotta},
  econ point/.style={circle, fill=SlateBlue!90!black, inner sep=1.8pt},
  econ highlight/.style={fill=Terracotta!20, opacity=0.5},
  econ label/.style={font=\small\sffamily, align=center},
  econ annotation/.style={font=\footnotesize\itshape, color=SlateBlue!80!black}
}


\begin{document}

\title{\Huge\bfseries\color{SlateBlue} Chapter 2: The Production Possibilities Frontier \& Gains from Trade}
\author{\large Principles of Microeconomics Curriculum}
\date{2026}
\maketitle

\thispagestyle{plain}

\begin{abstract}
This chapter provides a rigorous, pedagogical treatment of Chapter 2: The Production Possibilities Frontier \& Gains from Trade, benchmarked to standard principles curriculum.
\end{abstract}

\tableofcontents
\vspace{1.5em}
\hrule
\vspace{1.5em}

\section{The Production Possibilities Frontier (PPF) Model}
The PPF illustrates the boundary of maximum output combinations achievable under fixed technology and full factor utilization.
\begin{equation}
MRT = -\frac{\Delta Y}{\Delta X} = \left|\frac{dY}{dX}\right| = \frac{\text{Quantity of Good } Y \text{ Forgone}}{\text{Quantity of Good } X \text{ Gained}}
\end{equation}

\section{Law of Increasing Opportunity Cost}
Concave curvature reflects imperfect factor adaptability. As resources shift toward one sector, inputs progressively less suited to that industry must be utilized, driving opportunity costs upward.

\section{Ricardian Comparative Advantage}
Comparative advantage is determined strictly by lower domestic opportunity cost.
\begin{coreprinciple}
\boxheader{Terracotta}{Terms of Trade Rule}
Mutually beneficial exchange requires the Terms of Trade ($TOT$) to settle strictly between both domestic opportunity cost ratios:
\begin{equation}
\left|\frac{\Delta Y}{\Delta X}\right|_{\text{Region A}} < TOT < \left|\frac{\Delta Y}{\Delta X}\right|_{\text{Region B}}
\end{equation}
Specialization and trade shift the Consumption Possibilities Frontier ($CPF$) strictly outside the domestic autarky PPF.
\end{coreprinciple}

\end{document}
